\documentclass[10pt,a4paper]{amsart}
\usepackage{amsmath,amssymb,amsthm,graphicx}
% Retained minimal preamble for the Stacks Project chapter more-algebra.tex (More on Algebra).
% The chapter file pinned-original-source.tex carries only \input{preamble} on its first line, so no
% preamble text lives inside the chapter itself.  The Stacks Project uses one single shared preamble for
% every chapter; the snapshot stages it identically as sources/stacks-chow/preamble.tex and
% sources/stacks-etale-cohomology/preamble.tex.  The block below is copied mechanically, line for line,
% from sources/stacks-chow/preamble.tex lines 14-18 (the Xy-pic setup) and 154-217 (the native theorem
% declarations together with the mathematical macros).  Nothing was added, removed or redefined; the
% wrapper class and the amsmath/amssymb/amsthm/graphicx base packages are supplied by the assembler.
% Chapter: more-algebra.tex -- More on Algebra.
% For commutative diagrams we use Xy-pic
\usepackage[all]{xy}

% We use 2cell for 2-commutative diagrams.
\xyoption{2cell}

% Theorem environments.
%
\theoremstyle{plain}
\newtheorem{theorem}[subsection]{Theorem}
\newtheorem{proposition}[subsection]{Proposition}
\newtheorem{lemma}[subsection]{Lemma}

\theoremstyle{definition}
\newtheorem{definition}[subsection]{Definition}
\newtheorem{example}[subsection]{Example}
\newtheorem{exercise}[subsection]{Exercise}
\newtheorem{situation}[subsection]{Situation}

\theoremstyle{remark}
\newtheorem{remark}[subsection]{Remark}
\newtheorem{remarks}[subsection]{Remarks}

\numberwithin{equation}{subsection}

% Macros
%
\def\lim{\mathop{\mathrm{lim}}\nolimits}
\def\colim{\mathop{\mathrm{colim}}\nolimits}
\def\Spec{\mathop{\mathrm{Spec}}}
\def\Hom{\mathop{\mathrm{Hom}}\nolimits}
\def\Ext{\mathop{\mathrm{Ext}}\nolimits}
\def\SheafHom{\mathop{\mathcal{H}\!\mathit{om}}\nolimits}
\def\SheafExt{\mathop{\mathcal{E}\!\mathit{xt}}\nolimits}
\def\Sch{\mathit{Sch}}
\def\Mor{\mathop{\mathrm{Mor}}\nolimits}
\def\Ob{\mathop{\mathrm{Ob}}\nolimits}
\def\Sh{\mathop{\mathit{Sh}}\nolimits}
\def\NL{\mathop{N\!L}\nolimits}
\def\CH{\mathop{\mathrm{CH}}\nolimits}
\def\proetale{{pro\text{-}\acute{e}tale}}
\def\etale{{\acute{e}tale}}
\def\QCoh{\mathit{QCoh}}
\def\Ker{\mathop{\mathrm{Ker}}}
\def\Im{\mathop{\mathrm{Im}}}
\def\Coker{\mathop{\mathrm{Coker}}}
\def\Coim{\mathop{\mathrm{Coim}}}

% Boxtimes
%
\DeclareMathSymbol{\boxtimes}{\mathbin}{AMSa}{"02}

%
% Macros for moduli stacks/spaces
%
\def\QCohstack{\mathcal{QC}\!\mathit{oh}}
\def\Cohstack{\mathcal{C}\!\mathit{oh}}
\def\Spacesstack{\mathcal{S}\!\mathit{paces}}
\def\Quotfunctor{\mathrm{Quot}}
\def\Hilbfunctor{\mathrm{Hilb}}
\def\Curvesstack{\mathcal{C}\!\mathit{urves}}
\def\Polarizedstack{\mathcal{P}\!\mathit{olarized}}
\def\Complexesstack{\mathcal{C}\!\mathit{omplexes}}
% \Pic is the operator that assigns to X its picard group, usage \Pic(X)
% \Picardstack_{X/B} denotes the Picard stack of X over B
% \Picardfunctor_{X/B} denotes the Picard functor of X over B
\def\Pic{\mathop{\mathrm{Pic}}\nolimits}
\def\Picardstack{\mathcal{P}\!\mathit{ic}}
\def\Picardfunctor{\mathrm{Pic}}
\def\Deformationcategory{\mathcal{D}\!\mathit{ef}}
\usepackage{cleveref}
\crefname{originallemma}{original lemma}{original lemmas}
\Crefname{originallemma}{Original lemma}{Original lemmas}
\makeatletter
\newlabel{more-algebra-lemma-points-of-fibre-product}{{0B7J}{source}{}{sourcepointer-1}{}}
\newlabel{more-algebra-lemma-points-of-fibre-product@cref}{{[originallemma][0][]0B7J}{[1][0][]source}{}{}{}}
\makeatother
\crefname{originallemma}{original lemma}{original lemmas}
\Crefname{originallemma}{Original lemma}{Original lemmas}
\makeatletter
\newlabel{situation-module-over-fibre-product}{{08KH}{source}{}{sourcepointer-2}{}}
\newlabel{situation-module-over-fibre-product@cref}{{[originallemma][0][]08KH}{[1][0][]source}{}{}{}}
\makeatother
\crefname{originallemma}{original lemma}{original lemmas}
\Crefname{originallemma}{Original lemma}{Original lemmas}
\makeatletter
\newlabel{topology-section-colimits}{{0B1W}{source}{}{sourcepointer-3}{}}
\newlabel{topology-section-colimits@cref}{{[originallemma][0][]0B1W}{[1][0][]source}{}{}{}}
\makeatother
\crefname{originallemma}{original lemma}{original lemmas}
\Crefname{originallemma}{Original lemma}{Original lemmas}
\makeatletter
\newlabel{lemma-diagram-localize}{{01Z8}{source}{}{sourcepointer-4}{}}
\newlabel{lemma-diagram-localize@cref}{{[originallemma][0][]01Z8}{[1][0][]source}{}{}{}}
\makeatother
\title{Points of the spectrum of a fibre product of rings}
\author{The Stacks Project Authors}
\date{stacks-project:more-algebra.tex:local-snapshot-20261012; GNU FDL 1.2 or later}
\begin{document}
\maketitle
\setcounter{subsection}{1}
\noindent\textit{Editorial note.} All delivered mathematical text is the author's own, copied from the pinned Stacks Project chapter file more-algebra.tex; no statement or proof step was rewritten, reordered or removed.\par\medskip
\noindent\textit{Editorial note.} The chapter file pinned-original-source.tex carries only \textbackslash{}input\{preamble\} on its first line, so no preamble text lives inside the chapter. The retained file sources/stacks-more-algebra/retained-preamble.tex is a mechanical, line-for-line copy of the project's single shared preamble as staged in the snapshot (Xy-pic setup plus theorem declarations and mathematical macros). Nothing was added, removed or redefined, and the excerpt serves only as preamble, never as a span coordinate.\par\medskip
\noindent\textit{Editorial note.} Both counted spans start at the \textbackslash{}begin line of their environment, so the statement and the author's proof are delivered whole; surrounding prose of the section is deliberately outside the counted spans and is carried as standard prerequisites instead.\par\medskip
\noindent\textit{Editorial note.} The citation fields (source\_id, source\_version, authors, source\_url, license.name, source\_title) are identical to those of every other card this lane has produced from this chapter, so the batch selector treats them as one source.\par\medskip
\noindent\textit{Allowed standard prerequisites (editorial).}\begin{itemize}
\item The Situation of a fibre product of rings, with B' = B x\_A A' and the two structure maps. situation environment at source lines 521-542, immediately before the counted statement. Fixes all the rings used in the statement and in the proof.
\item Spectrum of a ring and the induced maps between spectra, including the description of fibres over a prime. standard, used throughout the proof at lines 553-604. Supplies the points whose description is the statement.
\item Prime ideals of a fibre product of rings and the correspondence with compatible pairs of primes. used in the two case discussions of the proof, lines 565-604. Provides the construction of the required primes and the verification of primeness.
\end{itemize}\medskip
\ifdefined\hypertarget\hypertarget{sourcepointer-1}{}\fi
\noindent\textit{Source reference.} Original lemma 0B7J: focused-seeds/AGP/sources/stacks-snapshot-20261011/tags.txt, source lines 14107--14107. The lemma's own canonical tag, quoted so that readers can resolve the label inside the delivered text; it points at the counted statement itself and carries no proof content.\par\medskip
\ifdefined\hypertarget\hypertarget{sourcepointer-2}{}\fi
\noindent\textit{Source reference.} Original lemma 08KH: focused-seeds/AGP/sources/stacks-snapshot-20261011/tags.txt, source lines 10885--10885. Navigation pointer to a statement the counted proof refers to but does not contain; nothing load-bearing is delegated to it.\par\medskip
\ifdefined\hypertarget\hypertarget{sourcepointer-3}{}\fi
\noindent\textit{Source reference.} Original lemma 0B1W: focused-seeds/AGP/sources/stacks-snapshot-20261011/tags.txt, source lines 13909--13909. Navigation pointer to a statement the counted proof refers to but does not contain; nothing load-bearing is delegated to it.\par\medskip
\ifdefined\hypertarget\hypertarget{sourcepointer-4}{}\fi
\noindent\textit{Source reference.} Original lemma 01Z8: focused-seeds/AGP/sources/stacks-snapshot-20261011/tags.txt, source lines 2583--2583. Navigation pointer to a statement the counted proof refers to but does not contain; nothing load-bearing is delegated to it.\par\medskip
\begin{lemma}
\label{lemma-points-of-fibre-product}
In Situation \ref{situation-module-over-fibre-product}
we have
$$
\Spec(B') = \Spec(B) \amalg_{\Spec(A)} \Spec(A')
$$
as topological spaces.
\end{lemma}

\begin{proof}
Since $B' = B \times_A A'$ we obtain a commutative square of
spectra, which induces a continuous map
$$
can : \Spec(B) \amalg_{\Spec(A)} \Spec(A') \longrightarrow \Spec(B')
$$
as the source is a pushout in the category
of topological spaces (which exists by
Topology, Section \ref{topology-section-colimits}).

\medskip\noindent
To show the map $can$ is surjective, let $\mathfrak q' \subset B'$ be a prime
ideal. If $I \subset \mathfrak q'$ (here and below we take the liberty
of considering $I$ as an ideal of $B'$ as well as an ideal of $A'$),
then $\mathfrak q'$ corresponds to a prime ideal of $B$ and is in the image.
If not, then pick $h \in I$, $h \not \in \mathfrak q'$.
In this case $B_h = A_h = 0$
and the ring map $B'_h \to A'_h$ is an isomorphism, see
Lemma \ref{lemma-diagram-localize}. Thus we see that
$\mathfrak q'$ corresponds to a unique prime ideal $\mathfrak p' \subset A'$
which does not contain $I$.

\medskip\noindent
Since $B' \to B$ is surjective, we see that $can$ is injective on
the summand $\Spec(B)$. We have seen above that $\Spec(A') \to \Spec(B')$
is injective on the complement of $V(I) \subset \Spec(A')$. Since
$V(I) \subset \Spec(A')$ is exactly the image of $\Spec(A) \to \Spec(A')$
a trivial set theoretic argument shows that $can$ is injective.

\medskip\noindent
To finish the proof we have to show that $can$ is open. To do this, observe
that an open of the pushout is of the form $V \amalg U'$ where
$V \subset \Spec(B)$ and $U' \subset \Spec(A')$ are opens whose inverse
images in $\Spec(A)$ agree. Let $v \in V$. We can find a $g \in B$
such that $v \in D(g) \subset V$. Let $f \in A$ be the image.
Pick $f' \in A'$ mapping to $f$. Then
$D(f') \cap U' \cap V(I) = D(f') \cap V(I)$.
Hence $V(I) \cap D(f')$ and $D(f') \cap (U')^c$ are disjoint closed
subsets of $D(f') = \Spec(A'_{f'})$. Write $(U')^c = V(J)$
for some ideal $J \subset A'$. Since
$A'_{f'} \to A'_{f'}/IA'_{f'} \times A'_{f'}/JA'_{f'}$
is surjective by the disjointness just shown, we can find an $a'' \in A'_{f'}$
mapping to $1$ in $A'_{f'}/IA'_{f'}$ and mapping to zero in
$A'_{f'}/JA'_{f'}$. Clearing denominators, we find an element
$a' \in J$ mapping to $f^n$ in $A$. Then $D(a'f') \subset U'$.
Let $h' = (g^{n + 1}, a'f') \in B'$.
Since $B'_{h'} = B_{g^{n + 1}} \times_{A_{f^{n + 1}}} A'_{a'f'}$ by a previously
cited lemma, we see that $D(h')$ pulls back to an open neighbourhood
of $v$ in the pushout, i.e., the image of $V \amalg U'$ contains
an open neighbourhood of the image of $v$. We omit the (easier) proof that
the same thing is true for $u' \in U'$ with $u' \not \in V(I)$.
\end{proof}

\end{document}
